%\documentclass{beamer}
%\documentclass[slidestop,usepdftitle=false]{beamer}
\documentclass[english,10pt,table]{beamer}
%\documentclass[english,10pt,table,handout]{beamer}

% Copyright 2007 by Till Tantau
%
% This file may be distributed and/or modified
%
% 1. under the LaTeX Project Public License and/or
% 2. under the GNU Public License.
%
% See the file doc/licenses/LICENSE for more details.


% Common packages

\usepackage[english]{babel}
\usepackage[utf8]{vietnam}
%\usepackage{times}
\usefonttheme[onlymath]{serif}
\usecolortheme{default}
\usepackage{booktabs}
\usepackage{mathpartir}
\usepackage{listings}
\usepackage{pbox}
\mprset{flushleft}
\mode<article>
{
  \usepackage{times}
  \usepackage{mathptmx}
  \usepackage[left=1.5cm,right=6cm,top=1.5cm,bottom=3cm]{geometry}
}

\usepackage{hyperref}
\usepackage{tikz}
\usetikzlibrary{arrows,backgrounds}
%\tikzstyle{mnode}=[circle, draw, fill=black, inner sep=0pt, minimum width=4pt]
\usepackage{colortbl}
%\usepackage{yfonts}
\usepackage{translator} % comment this, if not available


% Common settings for all lectures in this course

\def\lecturename{Cấu trúc rời rạc cho KHMT}

\title{\insertlecture}

\author{N.T.Toàn, P.N.Minh, L.H.Anh, T.Đ.G.Hải, H.Đ.Thịnh, N.P.Thịnh}

\institute
{
  Khoa Khoa học và Kỹ thuật Máy tính\\
  Trường Đại học Bách Khoa - ĐHQG TPHCM \\
	
}

\subject{Lecturer \lecturename}




% Beamer version theme settings

\useoutertheme[height=0pt,width=2cm,right]{sidebar}
\usecolortheme{rose,sidebartab}
\useinnertheme{circles}
\usefonttheme[only large]{structurebold}

\setbeamercolor{sidebar right}{bg=black!15}
\setbeamercolor{structure}{fg=blue}
\setbeamercolor{author}{parent=structure}

\setbeamerfont{title}{series=\normalfont,size=\LARGE}
\setbeamerfont{title in sidebar}{series=\bfseries}
\setbeamerfont{author in sidebar}{series=\bfseries}
\setbeamerfont*{item}{series=}
\setbeamerfont{frametitle}{size=}
\setbeamerfont{block title}{size=\small}
\setbeamerfont{subtitle}{size=\normalsize,series=\normalfont}

\setbeamertemplate{navigation symbols}{}
\setbeamertemplate{bibliography item}[book]
\setbeamertemplate{sidebar right}
{
  {\usebeamerfont{title in sidebar}%
    \vskip1.5em%
    \hskip3pt%
    \usebeamercolor[fg]{title in sidebar}%
    \insertshorttitle[width=2cm,center,respectlinebreaks]\par%
    \vskip1.25em%
  }%
  {%
    \hskip3pt%
    \usebeamercolor[fg]{author in sidebar}%
    \usebeamerfont{author in sidebar}%
    \insertshortauthor[width=2cm,center,respectlinebreaks]\par%
    \vskip1.25em%
  }%
  \hbox to2cm{\hss\insertlogo\hss}
  \vskip1.25em%
  \insertverticalnavigation{2cm}%
  \vfill
  \hbox to 2cm{\hfill\usebeamerfont{subsection in
      sidebar}\strut\usebeamercolor[fg]{subsection in
      sidebar}\insertshortlecture.\insertframenumber\hskip5pt}%
  \vskip3pt%
}%

\setbeamertemplate{title page}
{
  \vbox{}
  \vskip1em
  {\normal ỨNG DỤNG THỐNG KÊ  \par}
  {\usebeamercolor[fg]{title}\usebeamerfont{title}\inserttitle\par}%
  \ifx\insertsubtitle\@empty%
  \else%
    \vskip0.25em%
    {\usebeamerfont{subtitle}\usebeamercolor[fg]{subtitle}\insertsubtitle\par}%
  \fi%     
  \vskip1em\par
  \emph{\lecturename}\ on \insertdate\par
  \vskip0pt plus1filll
  \leftskip=0pt plus1fill\insertauthor\par
  \insertinstitute\vskip1em
}

\logo{\includegraphics[width=1.5cm]{hcmut.png}}



% Article version layout settings

\mode<article>

\makeatletter
\def\@listI{\leftmargin\leftmargini
  \parsep 0pt
  \topsep 5\p@   \@plus3\p@ \@minus5\p@
  \itemsep0pt}
\let\@listi=\@listI


\setbeamertemplate{frametitle}{\paragraph*{\insertframetitle\
    \ \small\insertframesubtitle}\ \par
}
\setbeamertemplate{frame end}{%
  \marginpar{\scriptsize\hbox to 1cm{\sffamily%
      \hfill\strut\insertshortlecture.\insertframenumber}\hrule height .2pt}}
\setlength{\marginparwidth}{1cm}
\setlength{\marginparsep}{4.5cm}

\def\@maketitle{\makechapter}

\def\makechapter{
  \newpage
  \null
  \vskip 2em%
  {%
    \parindent=0pt
    \raggedright
    \sffamily
    \vskip8pt
    {\fontsize{36pt}{36pt}\selectfont Chapter \insertshortlecture \par\vskip2pt}
    {\fontsize{24pt}{28pt}\selectfont \color{blue!50!black} \insertlecture\par\vskip4pt}
    {\Large\selectfont \color{blue!50!black} \insertsubtitle\par}
    \vskip10pt

  }
  \par
  \vskip 1.5em%
}

\let\origstartsection=\@startsection
\def\@startsection#1#2#3#4#5#6{%
  \origstartsection{#1}{#2}{#3}{#4}{#5}{#6\normalfont\sffamily\color{blue!50!black}\selectfont}}

\makeatother

\mode
<all>




% Typesetting Listings

\usepackage{listings}
\lstset{language=Java}

\alt<presentation>
{\lstset{%
  basicstyle=\footnotesize\ttfamily,
  commentstyle=\slshape\color{green!50!black},
  keywordstyle=\bfseries\color{blue!50!black},
  identifierstyle=\color{blue},
  stringstyle=\color{orange},
  escapechar=\#,
  emphstyle=\color{red}}
}
{
  \lstset{%
    basicstyle=\ttfamily,
    keywordstyle=\bfseries,
    commentstyle=\itshape,
    escapechar=\#,
    emphstyle=\bfseries\color{red}
  }
}



% Common theorem-like environments

\theoremstyle{example}
\newtheorem{exercise}[theorem]{\translate{Exercise}}


% New useful definitions:

\newbox\mytempbox
\newdimen\mytempdimen

\newcommand\includegraphicscopyright[3][]{%
  \leavevmode\vbox{\vskip3pt\raggedright\setbox\mytempbox=\hbox{\includegraphics[#1]{#2}}%
    \mytempdimen=\wd\mytempbox\box\mytempbox\par\vskip1pt%
    \fontsize{3}{3.5}\selectfont{\color{black!25}{\vbox{\hsize=\mytempdimen#3}}}\vskip3pt%
}}

\newenvironment{colortabular}[1]{\medskip\rowcolors[]{1}{blue!20}{blue!10}\tabular{#1}\rowcolor{blue!40}}{\endtabular\medskip}

\def\equad{\leavevmode\hbox{}\quad}

\newenvironment{greencolortabular}[1]
{\medskip\rowcolors[]{1}{green!50!black!20}{green!50!black!10}%
  \tabular{#1}\rowcolor{green!50!black!40}}%
{\endtabular\medskip}




\lecture[0]{Khảo sát kết quả của bài tập online cho phép nộp bài nhiều lần}{lecture-text}

\usepackage{pifont}
% Symbol definitions for these lists
\newcommand{\DingListSymbolA}{43}
\newcommand{\DingListSymbolB}{243}
\newcommand{\DingListSymbolC}{224}
\newcommand{\DingListSymbolD}{219}

% Boxed equation
\definecolor{LightYellow}{rgb}{1.,1.,.9}
\definecolor{LightRed}{rgb}{1.,.6,.6}


%%ensembles de nombres
\def\NP{$\mathcal{NP}$}
\def\N{\mathbb{N}}
\def\Z{\mathbb{Z}}
\def\R{\mathbb{R}}
\def\Q{\mathbb{Q}}

%\date[]{~~}
%%%%%%%%%%%%%%%%%%%%%%%%%%%%%%%%%%%%%%%%%%%%%%%%%%%%%%%%%%%%%%%%%%%%%
%%%%%%%%%%%%%%%%%%%%%%%%%%%%%%%%%%%%%%%%%%%%%%%%%%%%%%%%%%%%%%%%%%%%%
\begin{document}
\frame{
\selectlanguage{vietnamese}
  \maketitle
}


%%%%%%%%%%%%%%%%%%%%%%%%%%%%%%%%%%%%%%%%%%%%%%%%%%%%%%%%%%%%%%%%%%%%%
%%%%%%%%%%%%%%%%%%%%%%%%%%%%%%%%%%%%%%%%%%%%%%%%%%%%%%%%%%%%%%%%%%%%%
%\section[Plan]{}
%\setcounter{tocdepth}{1}
\frame{ \tableofcontents}
%\setcounter{tocdepth}{5}
% to display left summary deeper and plan slide juste display section: add command \setcounter{tocdepth}{1} and then \setcounter{tocdepth}{10}  recompile twice or more again 

%%%%%%%%%%%%%%%%%%%%%%%%%%%%%%%%%%%%%%%%%%%%%%%%%%%%%%%%%%%%%%%%%%%%%
%%%%%%%%%%%%%%%%%%%%%%%%%%%%%%%%%%%%%%%%%%%%%%%%%%%%%%%%%%%%%%%%%%%%%
\section{File 1}
\subsection{Bài 1}
\frame
{
  \frametitle{Bài 1: Xác định số lượng sinh viên}	
\begin{block}{Đâu nào là một mệnh đề}
Xác định giá trị chân lý của mệnh đề đó và tìm phủ định của nó.
\begin{enumerate}[a]
	\item Miami là thủ phủ của bang Florida.
	\item 1+1=2.
	\item Hãy tìm giá trị của $x$ nếu biết $x+2=5$
	\item Đây là đường một chiều.
	\item Bây giờ là mấy giờ?
	\item Đất đỏ bazane trồng cây rất tốt.
	\item $x+y=y+z$ nếu $x=z$.
	\item Hôm nay là thứ năm
	\item Không có ô nhiễm ở Bảo Lộc, nhưng có nhiều ô nhiễm ở Đà Lạt.
	\item Mùa hè ở thành phố Seville thì nóng và nắng.
\end{enumerate}
\end{block}
}

\frame
{
  \frametitle{Bài 1: Xác định số lượng sinh viên}	
\begin{block}{}
Chứng minh các mệnh đề sau đây tương đương: 
\begin{enumerate}[a]
	\item $\lnot (p \leftrightarrow q) $ và $\lnot p \leftrightarrow q$ 
	\item $(p \to q) \land (p \to r)$ và $p \to (q \land r)$  
	\item $(p \to r) \land (q \to r)$ và $(p \lor q) \to r$
	\item $(p \to q) \lor (p \to r)$ và $p \to (q \lor r)$ 
	\item $\lnot p \to (q \to r)$ và $q \to (p \lor r)$ 
	\item $p \leftrightarrow q$ và $(p \to q) \land (q \to p)$ 
\end{enumerate}
\end{block}
}

\frame
{
  \frametitle{Proving methods}	
\begin{block}{}
Gọi $P$, $Q$ là các mệnh đề: 
\begin{itemize}
	\item $P$: ``Minh giỏi Toán''
	\item $Q$: ``Minh yếu Anh văn''
\end{itemize}
Hãy viết lại các mệnh đề sau dưới dạng hình thức trong đó sử dụng các phép nối. \\
(Giả sử đảo nghĩa của học 'giỏi' là học 'yếu'.)
\begin{enumerate}[a]
\item Minh giỏi Toán nhưng yếu Anh văn
\item Minh yếu cả Toán lẫn Anh văn
\item Minh giỏi Toán hay Minh vừa giỏi Anh văn nhưng vừa yếu Toán
\item Nếu Minh giỏi Toán thì Minh giỏi Anh văn
\item Minh giỏi Toán và Anh văn hay Minh yếu Toán nhưng giỏi Anh
\end{enumerate}

\end{block}
}
%%%%%%%%%%%%%%%%%%%%%%%%%%%%%%%
\subsection{Bài 2}
\frame
{
  \frametitle{Predicate logic}	
\begin{block}{}
Cho vị từ $N(x)$ ``$x$ đã từng đi chơi Đà Lạt'' với tập vũ trụ là toàn bộ sinh viên trong lớp Cấu trúc rời rạc. 
Hãy phát biểu các vị từ sau:
\begin{enumerate}[a]
	\item $\exists x N(x)$ 
	\item $\forall x N(x)$ 
	\item $\neg \exists x N(x)$ 
 \item $\exists x \neg N(x) $ 
	\item $\neg \forall x N(x)$ 
	\item $\forall x \neg N(x)$
\end{enumerate}
\end{block}
}

%%%%%%%%%%%%%%%%%%%%%%%%%%%%%%%
\subsection{Bài 4}
\frame
{
  \frametitle{Proving methods}	
\begin{block}{}
Show that if $n$ is an integer and $n^3 + 2015$ is odd, then $n$ is even using
\begin{enumerate}[a]
\item a proof by contraposition.
\item a proof by contradiction.
\end{enumerate}
\end{block}
}


%%%%%%%%%%%%%%%%%%%%%%%%%%%%%%%%%%%%%%%%%%%%%%%%%%%%%%%%%%%%%%%%%%%%%
%%%%%%%%%%%%%%%%%%%%%%%%%%%%%%%%%%%%%%%%%%%%%%%%%%%%%%%%%%%%%%%%%%%%%
\subsection{Bài 5}
\frame{}
\subsection{Bài 7}
\frame{}
\subsection{Bài 9}
\frame{}
\end{document}

